% Blueprint note: Fixed upstream source versions: \cite{KIPSource5aa845241533,KIPSourcebbac1fc05843}. Interface scope notes are retained as native hgraph node comments; historical migration records are local archives.

% Curated mathematical exposition; compiled dependencies and status are supplied by the publisher.

\section{Archived generators and bounded coordinates}

\begin{definition}[Archived labels and their bidegrees]
\label{sem:Computation:archive-labels}
\level{coarse}
From \cite{LinProgram}, archived sphere generators.\par
The selected labels $h_i,g,e_0,\Delta h_1g,\Delta h_6g,Md_0,[B_4],b$ and $x_{n,s,j}$ carry their exact archived index, name and Adams bidegree $(s,t)$. A generator catalogue records every generator and its degree; a source monomial is formed from its unit, generators and degree-preserving products.
\lean{KIP126.Computation.Atom,KIP126.Computation.Atom.record,KIP126.Computation.Atom.degree,KIP126.Computation.GeneratorCatalogue,KIP126.Computation.GeneratorCatalogue.degree,KIP126.Computation.SourceMonomial}
\end{definition}

\begin{definition}[The source presentation and complete archive]
\label{sem:Computation:source-presentation}
\level{coarse}
\uses{sem:Computation:archive-labels}
A source presentation consists of bigraded $\mathbb F_2$-modules, specified bilinear products, a unit element and the named atoms. A sphere archive specifies its complete generator list, the exact identification of every atom, and every basis monomial in each bidegree. Evaluation of a monomial uses these same products and generator values.
\lean{KIP126.Computation.ProgramData,KIP126.Computation.SphereArchive,KIP126.Computation.SourceMonomial.value}
\end{definition}

\begin{definition}[The fixed sphere archive]
\label{sem:Computation:fixed-archive}
\level{coarse}
\uses{sem:Computation:source-presentation}
From \cite{LinProgram}, archived sphere basis.\par
Fix the project generator catalogue, atom indices, basis counts and basis monomials. The atom index determines the archived generator and its prescribed degree.
\lean{KIP126.Computation.Standard.generators,KIP126.Computation.Standard.atomIndex,KIP126.Computation.Standard.atom_record,KIP126.Computation.Standard.basisCount,KIP126.Computation.Standard.basisMonomial,KIP126.Computation.Standard.archive,KIP126.Computation.SphereArchive.atom_degree}
\end{definition}

\begin{definition}[The fixed source components and products]
\label{sem:Computation:source-components}
\level{coarse}
\uses{sem:Computation:fixed-archive}
The fixed source object has components $P^{s,t}$ over $\mathbb F_2$, bilinear maps $P^{s,t}\times P^{s\prime,t\prime}\to P^{s+s\prime,t+t\prime}$, a unit in degree $(0,0)$ and its prescribed atom in each named degree. For each atom, its selected generator has the prescribed bidegree and value.
\lean{KIP126.Computation.Standard.sourceComponent,KIP126.Computation.Standard.sourceProduct,KIP126.Computation.Standard.sourceUnit,KIP126.Computation.Standard.sourceAtom,KIP126.Computation.Standard.sourceData,KIP126.Computation.Standard.sourceGenerator,KIP126.Computation.Standard.sourceAtom_agrees}
\end{definition}

\begin{definition}[Complete source coordinates]
\label{sem:Computation:source-basis}
\level{coarse}
\uses{sem:Computation:source-components,sem:Computation:fixed-archive}
For $s,t\in\mathbb N$ and $t\leq261$, the source coordinates are a linear equivalence
\[P^{s,t}\cong\bigoplus_{i<\operatorname{basisCount}(s,t)}\mathbb F_2.\]
The inverse of each coordinate unit vector is the value of the corresponding archived monomial. In particular, the archived monomials form a basis of $P^{s,t}$.
\lean{KIP126.Computation.Standard.sourceCoordinates,KIP126.Computation.Standard.sourceBasis_values}
\end{definition}

\section{Comparison with the actual sphere}

\begin{definition}[The bounded comparison]
\label{sem:Computation:bounded-comparison}
\level{coarse}
\uses{sem:Computation:source-components,sem:AdamsSynthetic:families,sem:AdamsSynthetic:standard-classes}
For $s,t\in\mathbb N$ and $t\leq261$, the fixed comparison is a linear equivalence from $P^{s,t}$ to the actual sphere $E_2^{s,t}$. A sphere presentation requires the comparison to preserve the specified unit and every product of total degree $t+t\prime\leq261$, and to identify the selected $h_i$ with the fixed standard classes.
\lean{KIP126.Computation.SpherePresentation,KIP126.Computation.Standard.comparison}
\end{definition}

\begin{theorem}[Products and the unit under comparison]
\label{sem:Computation:multiplicative-comparison}
\level{coarse}
\uses{sem:Computation:bounded-comparison,frontier:AdamsSynthetic:f05463b82afe}
If $t+t\prime\leq261$, then
\[c(xy)=c(x)c(y),\qquad c(1)=1.\]
The comparison is $\mathbb F_2$-linear.
\lean{KIP126.Computation.Standard.comparison_product,KIP126.Computation.Standard.comparison_unit,KIP126.Computation.Standard.comparisonF2,KIP126.Computation.SpherePresentation.comparisonF2,KIP126.Computation.SpherePresentation.comparison_mul}
\end{theorem}

\begin{theorem}[Identification of the standard atoms]
\label{sem:Computation:standard-atoms}
\level{coarse}
\uses{sem:Computation:bounded-comparison,sem:AdamsSynthetic:standard-classes}
For $i=0,1,2,4,5,6$, the source atom $h_i$ is sent to the designated standard class $h_i\in E_2^{1,2^i}$. For every archived atom $a$, define its page class by $c(a)$.
\lean{KIP126.Computation.Standard.atom,KIP126.Computation.Standard.comparison_h0,KIP126.Computation.Standard.comparison_h1,KIP126.Computation.Standard.comparison_h2,KIP126.Computation.Standard.comparison_h4,KIP126.Computation.Standard.comparison_h5,KIP126.Computation.Standard.comparison_h6,KIP126.Computation.Standard.atom_h0,KIP126.Computation.Standard.atom_h1,KIP126.Computation.Standard.atom_h2,KIP126.Computation.Standard.atom_h4,KIP126.Computation.Standard.atom_h5,KIP126.Computation.Standard.atom_h6}
\end{theorem}

\begin{theorem}[Coordinates on the actual page]
\label{sem:Computation:actual-coordinates}
\level{coarse}
\uses{sem:Computation:source-basis,sem:Computation:bounded-comparison}
For $t\leq261$, composition of the inverse comparison with the source coordinates gives
\[E_2^{s,t}\cong\bigoplus_{i<\operatorname{basisCount}(s,t)}\mathbb F_2.\]
Its inverse sends each unit vector to the image of the same archived monomial. Pulling that vector back by the comparison recovers the original source monomial.
\lean{KIP126.Computation.Standard.coordinates,KIP126.Computation.Standard.coordinates_basis,KIP126.Computation.Standard.comparison_symm_coordinates_basis,KIP126.Computation.SpherePresentation.coordinates,KIP126.Computation.SpherePresentation.coordinates_basis,KIP126.Computation.SpherePresentation.comparison_symm_coordinates_basis}
\end{theorem}

\begin{theorem}[Calibration of the Kervaire class]
\label{sem:Computation:h6-square-calibration}
\level{coarse}
\uses{sem:Computation:standard-atoms,sem:Computation:multiplicative-comparison}
The source product of the archived $h_6$ with itself maps to the designated actual class $h_6^2\in E_2^{2,128}$.
\lean{KIP126.Computation.Standard.h6_square,KIP126.Computation.SpherePresentation.h6_square}
\end{theorem}

\section{Named E2 classes near the Kervaire stem}

\begin{definition}[Actual products and powers]
\label{sem:Computation:products-and-powers}
\level{coarse}
\uses{sem:Computation:standard-atoms,sem:Computation:multiplicative-comparison}
For $x\in E_2^{s,t}$ and $y\in E_2^{s\prime,t\prime}$, their product belongs to $E_2^{s+s\prime,t+t\prime}$. The classes $h_0^2$, $h_0^6$ and $h_5^2$ are the corresponding iterated products.
\lean{KIP126.Computation.Standard.product,KIP126.Computation.Standard.h0Square,KIP126.Computation.Standard.h0Six,KIP126.Computation.Standard.h5Square}
\end{definition}

\begin{definition}[The classes W, T and U]
\label{sem:Computation:w-t-u}
\level{coarse}
\uses{sem:Computation:standard-atoms,sem:Computation:products-and-powers}
Define
\[W=x_{126,8,4}+x_{126,8},\quad T=h_1h_4x_{109,12},\quad U=h_0^2x_{124,8}.\]
Their Adams bidegrees are $(8,134)$, $(14,139)$ and $(10,134)$, respectively.
\lean{KIP126.Computation.Standard.W,KIP126.Computation.Standard.T,KIP126.Computation.Standard.U}
\end{definition}

\begin{definition}[The classes V, X and Y]
\label{sem:Computation:v-x-y}
\level{coarse}
\uses{sem:Computation:standard-atoms,sem:Computation:products-and-powers}
Define
\[V=x_{123,9}+h_0x_{123,8},\quad X=h_1x_{121,7},\quad Y=h_0^2x_{125,9,2}.\]
Their bidegrees are $(9,132)$, $(8,130)$ and $(11,136)$. The class $h_1V+U$ is a specified element of bidegree $(10,134)$.
\lean{KIP126.Computation.Standard.V,KIP126.Computation.Standard.X,KIP126.Computation.Standard.Y,KIP126.Computation.Standard.d2Target}
\end{definition}

\begin{definition}[The classes P, Q and the correction term]
\label{sem:Computation:p-q-correction}
\level{coarse}
\uses{sem:Computation:standard-atoms,sem:Computation:products-and-powers}
Define $P=h_6Md_0\in E_2^{11,133}$, $Q=h_5x_{91,11}\in E_2^{12,134}$ and $C=e_0\Delta h_6g\in E_2^{13,137}$. Also set $H=g^4\Delta h_1g\in E_2^{25,150}$.
\lean{KIP126.Computation.Standard.P,KIP126.Computation.Standard.Q,KIP126.Computation.Standard.correction,KIP126.Computation.Standard.highClass}
\end{definition}

\begin{theorem}[The h6-square line]
\label{sem:Computation:square-line}
\level{coarse}
\uses{sem:Computation:actual-coordinates,sem:Computation:h6-square-calibration}
From \cite{LinProgram}, archived sphere $E_2$ basis.\par
The class $h_6^2$ is nonzero and every $x\in E_2^{2,128}$ satisfies $x=0$ or $x=h_6^2$.
\lean{KIP126.Computation.Standard.square_nonzero,KIP126.Computation.Standard.square_exhaustive}
\end{theorem}

\section{Products, divisibility and E2 vanishing}

\begin{definition}[The source label for the seventh differential]
\label{sem:E2Computation:d7-label}
\level{coarse}
\uses{sem:Computation:standard-atoms,sem:Computation:products-and-powers}
From \cite{LinProgram}, archived sphere coordinates.\par

Define the second-page class
\[A=x_{123,11,2}+x_{123,11}+h_0h_6[B_4]\in E_2^{11,134}.\]
Its later differential is stated in \ref{sem:Computation:d7-source}.
\lean{KIP126.Computation.Standard.d7Source}
\end{definition}

\begin{definition}[The two target labels in bidegree (9,134)]
\label{sem:E2Computation:d3-labels}
\level{coarse}
\uses{sem:Computation:standard-atoms,sem:Computation:products-and-powers}
From \cite{LinProgram}, archived sphere coordinates.\par

Define
\[B=h_5x_{94,8},\qquad B'=h_5x_{94,8}+h_6b\in E_2^{9,134}.\]
The differential alternative involving them is stated in
\ref{sem:Computation:d3-candidates}.
\lean{KIP126.Computation.Standard.d3Candidate,KIP126.Computation.Standard.d3OtherCandidate}
\end{definition}

\begin{theorem}[Product relations for the named classes]
\label{sem:Computation:products-vanishing}
\level{coarse}
\uses{sem:Computation:v-x-y,sem:Computation:p-q-correction,sem:Computation:named-basis-identities,sem:Computation:theta-product-certificates}
From \cite{LinProgram}, archived product relations.\par
The actual sphere products satisfy
\[h_5^2b=0,\qquad Xh_2=0,\qquad h_1C=0.\]
Furthermore $Ph_2$ and $Qh_2$ are boundaries by page $2$, in the sense of Definition~\ref{sem:SpectralSequence:cycles-boundaries}.
\lean{KIP126.Computation.Standard.h5Square_b_zero,KIP126.Computation.Standard.x_h2_zero,KIP126.Computation.Standard.h1_correction_zero}
\end{theorem}

\begin{theorem}[Nondivisibility of T and Y]
\label{sem:Computation:nondivisibility}
\level{coarse}
\uses{sem:Computation:w-t-u,sem:Computation:v-x-y,sem:Computation:actual-coordinates,sem:Computation:obstruction-certificates}
From \cite{LinProgram}, archived product calculation.\par
There is no $a\in E_2^{13,138}$ with $h_0a=T$, no $a\in E_2^{13,135}$ with $h_2a=T$, and no $a\in E_2^{10,132}$ with $h_2a=Y$.
\lean{KIP126.Computation.Standard.t_not_h0_multiple,KIP126.Computation.Standard.t_not_h2_multiple,KIP126.Computation.Standard.y_not_h2_multiple}
\end{theorem}

\begin{theorem}[The required vanishing components]
\label{sem:Computation:vanishing-regions}
\level{coarse}
\uses{sem:Computation:actual-coordinates}
From \cite{LinProgram}, archived sphere $E_2$ dimensions.\par
The actual components $E_2^{s,s+125}$ for $0\leq s\leq4$ are zero. Here zero means that the corresponding additive module is a subsingleton.
\lean{KIP126.Computation.Standard.e2_stem125_low}
\end{theorem}

\section{The nu-cofiber and detector E2 comparisons}

\begin{definition}[The actual nu-cofiber and its labels]
\label{sem:Computation:nu-cofiber}
\level{coarse}
\uses{frontier:AdamsSynthetic:0895d4999173,sem:Computation:standard-atoms,sem:AdamsSynthetic:families}
Let $C\nu$ be the cofiber of the fixed map $\nu$. Its Adams spectral sequence is the one assigned to this actual cofiber. Write $i_*$ for the bottom-cell page map and $a$ for the selected top-cell label in bidegree $(8,134)$. Define
\[A_{C\nu}=a+i_*x_{126,8,2}+i_*x_{126,8}.\]
\lean{KIP126.Computation.Standard.nuCofiberSequence,KIP126.Computation.Standard.bottom,KIP126.Computation.Standard.topCell,KIP126.Computation.Standard.cofiberSource}
\end{definition}

\begin{theorem}[Identification of the top-cell image]
\label{sem:Computation:top-cell-comparison}
\level{coarse}
\uses{sem:Computation:nu-cofiber,sem:Computation:v-x-y,sem:AdamsSynthetic:selected-convergence}
From \cite{LinProgram}, archived nu-cofiber comparison.\par
The image of $a$ under the actual cofiber connecting arrow desuspends, through the prescribed four sphere suspensions and shift coherence maps, to $X=h_1x_{121,7}$. The comparison supplies the intermediate labels in degrees $(8,133)$, $(8,132)$ and $(8,131)$.
\lean{KIP126.Computation.Standard.top_comparison}
\end{theorem}

\begin{theorem}[Agreement with the literature labels]
\label{sem:Computation:literature-labels}
\level{coarse}
\uses{sem:Computation:standard-atoms,sem:Literature:kappa-detection,sem:Literature:w-detection}
The images of the archived classes $g$ and $\Delta h_1g$ under $c$ equal the respective sphere classes of the same names in \cite{LWX126}.
\lean{KIP126.Computation.Standard.g_agrees,KIP126.Computation.Standard.delta_agrees}
\end{theorem}

\begin{definition}[The fixed detector comparison]
\label{sem:Computation:detector-comparison}
\level{coarse}
\uses{sem:Computation:archive-labels,sem:Literature:tmf-unit,sem:Literature:tmf-d-three}
The archived detector has bigraded $\mathbb F_2$-components, selected elements $v_2^{16}$ and $B_{19,114}$ in degrees $(16,112)$ and $(19,114)$, and a comparison with the actual detector $E_2$ for $t\leq261$. The two selected elements agree, under that comparison, with the fixed literature labels transported by the prescribed detector isomorphism.
\lean{KIP126.Computation.Standard.detectorComponent,KIP126.Computation.Standard.detectorV2Sixteen,KIP126.Computation.Standard.detectorBetaGFour,KIP126.Computation.Standard.detectorComparison,KIP126.Computation.Standard.detector_v2_agrees,KIP126.Computation.Standard.detector_beta_agrees}
\end{definition}

\section{Finite coordinates and calibrated polynomial values}

\begin{definition}[Finite table coordinates]
\label{sem:Computation:table-coordinate-syntax}
\level{coarse}
From \cite{LinProgram}, archived finite tables.\par
A table row specifies the sphere or $C\nu$, an Adams bidegree and an increasing list of coordinate indices. A row is coordinate-valid if its degree occurs in the table and its coordinate indices are strictly increasing and less than the tabulated dimension.
\lean{KIP126.Computation.Table.Raw.Spectrum,KIP126.Computation.Table.Raw.Degree,KIP126.Computation.Table.Raw.degrees,KIP126.Computation.Table.Raw.strictlyIncreasing,KIP126.Computation.Table.Raw.coordinatesValid}
\end{definition}

\begin{definition}[Realization and decoding of a table row]
\label{sem:Computation:table-realization}
\level{coarse}
\uses{sem:Computation:table-coordinate-syntax,sem:Computation:actual-coordinates,sem:Computation:nu-cofiber}
For $o=S^0$ or $C\nu$, write $b_{o;s,t,i}\in E_2^{s,t}(o)$ for the specified basis vectors. For the sphere, $t\leq261$ and $i<\operatorname{basisCount}(s,t)$, these are the inverse coordinate images of the unit vectors; set $b_{S^0;s,t,i}=0$ when either bound fails. A coordinate-valid list $I$ represents
\[\sum_{i\in I}b_{o;s,t,i}.\]
An invalid coordinate list represents no element. The basis is complete if there is a linear coordinate equivalence sending every $b_{o;s,t,i}$ to the corresponding unit vector.
\lean{KIP126.Computation.Table.sphereBasis,KIP126.Computation.Table.cofiberBasis,KIP126.Computation.Table.basis,KIP126.Computation.Table.decode,KIP126.Computation.Table.BasisCorrect}
\end{definition}

\begin{theorem}[Correctness of the finite table bases]
\label{sem:Computation:table-bases}
\level{coarse}
\uses{sem:Computation:table-realization,sem:Computation:fixed-archive}
For every degree listed in the finite archive, its specified basis is complete. On the sphere, its length equals the fixed archive basis count; on $C\nu$, correctness refers to the prescribed cofiber basis. For $t\leq261$, the complete sphere archive count equals the tabulated dimension in degree $(s,t)$.
\lean{KIP126.Computation.Table.basis_count,KIP126.Computation.Table.cofiber_basis_correct,KIP126.Computation.Table.basis_correct,KIP126.Computation.Table.archive_dimension}
\end{theorem}

\begin{definition}[Polynomial expressions and calibrated basis values]
\label{sem:Computation:polynomial-values}
\level{coarse}
\uses{sem:Computation:table-realization,sem:Computation:standard-atoms,sem:Computation:multiplicative-comparison}
Let $G_{j;s,t}$ be the image in $E_2^{s,t}$ of archived generator $j$ when $t\leq261$ and $j$ has bidegree $(s,t)$, and set $G_{j;s,t}=0$ otherwise. Expressions are formed from $1$, the named atoms, the elements $G_{j;s,t}$ and products. Their values are computed in the sphere $E_2$-page. For each basis-value record $(e,i)$, the $i$th basis vector in the bidegree of $e$ equals the value of $e$.
\lean{KIP126.Computation.Table.Expression,KIP126.Computation.Table.Expression.degree,KIP126.Computation.Table.Expression.value,KIP126.Computation.Table.generatorImage,KIP126.Computation.Table.BasisValue,KIP126.Computation.Table.basisValues,KIP126.Computation.Table.basis_value}
\end{definition}

\begin{theorem}[The named table basis elements]
\label{sem:Computation:named-basis-identities}
\level{coarse}
\uses{sem:Computation:polynomial-values,sem:Computation:w-t-u,sem:Computation:p-q-correction}
The sphere basis vectors satisfy
\[b_{S^0;11,133,1}=P,\qquad b_{S^0;12,134,0}=Q,\qquad b_{S^0;14,139,1}=T,\]
\[b_{S^0;10,134,4}=U,\qquad b_{S^0;13,137,1}=C,\qquad b_{S^0;13,137,2}=h_4x_{109,12}.\]
If $x\in E_2^{s,t}$, $y\in E_2^{s\prime,t\prime}$ and $t+t\prime\leq261$, then $xy=yx$. Moreover, $h_0h_1=h_1h_2=0$.
\lean{KIP126.Computation.Table.basis_P,KIP126.Computation.Table.basis_Q,KIP126.Computation.Table.basis_T,KIP126.Computation.Table.basis_U,KIP126.Computation.Table.basis_correction,KIP126.Computation.Table.basis_h4_x109,KIP126.Computation.Table.product_comm,KIP126.Computation.Table.h0_h1_zero,KIP126.Computation.Table.h1_h2_zero}
\end{theorem}

\section{Primitive relation certificates}

\begin{theorem}[The primitive obstruction certificates]
\label{sem:Computation:obstruction-certificates}
\level{coarse}
\uses{sem:Computation:polynomial-values,sem:Computation:v-x-y}
From \cite{LinProgram}, archived relation certificates.\par
The following relations hold in the sphere $E_2$-page:
\[h_0G_{24;5,36}=0,\qquad (h_0G_{3;1,8})G_{335;12,130}=h_0^2G_{425;12,137},\]
\[h_0^5x_{125,9,2}=0,\qquad h_0^6x_{125,8}=0,\qquad h_2G_{407;13,135}=0,\]
\[h_2Md_0=h_0G_{89;10,72},\qquad h_0^3G_{89;10,72}=0.\]
% Blueprint note: These seven identities are the primitive polynomial certificates; no later homotopy conclusion is included.
\lean{KIP126.Computation.Table.h0_obstruction_rel0,KIP126.Computation.Table.h0_obstruction_rel1,KIP126.Computation.Table.h0_obstruction_rel2,KIP126.Computation.Table.h0_obstruction_rel3,KIP126.Computation.Table.h2_obstruction_rel0,KIP126.Computation.Table.h2_obstruction_rel1,KIP126.Computation.Table.h2_obstruction_rel2}
\end{theorem}

\begin{theorem}[Primitive theta product relations]
\label{sem:Computation:theta-product-certificates}
\level{coarse}
\uses{sem:Computation:polynomial-values}
From \cite{LinProgram}, archived theta relation certificates.\par
The following products vanish:
\[h_5G_{84;7,70}=0,\qquad h_5G_{85;7,70}=0,\qquad h_5h_6=0.\]
% Blueprint note: These identities are inputs to later indeterminacy arguments and do not themselves compute a Massey product.
\lean{KIP126.Computation.Table.theta_h5_generator84_zero,KIP126.Computation.Table.theta_h5_generator85_zero,KIP126.Computation.Table.theta_h5_h6_zero}
\end{theorem}



\section{The unbounded second-line basis}
\begin{theorem}[Exhaustive standard products on the Adams second line]
\label{sem:E2Computation:two-line-exhaustive}
\level{coarse}
From \cite{Adams60}, Theorem 2.5.1(2), p.~45.\par

For every integer $t$ and nonzero $x\in E_2^{2,t}(S^0)$, there are natural numbers $i\leq j$ and a bidegree equation
\[
 (1,2^i)+(1,2^j)=(2,t)
\]
such that $x$ equals the corresponding transport of $H_{i,j}$. There is no upper bound on $t$, $i$ or $j$.
\lean{KIP126.E2Computation.two_line_exhaustive}
% Blueprint note: Original source checked directly: \cite{Adams60}, Theorem 2.5.1(2), pp.~44--45, gives the second-line basis h_j h_i for j>=i>=0 and j!=i+1. \cite{LinExt4}, Theorem 1.2, p.~460, recalls this earlier result and is not its original source. Distinct ordered power-of-two pairs have distinct sums, so each nonzero homogeneous element over F2 is a single basis product. The theorem is adapted to integer grading, actual page products and their degree transports. Identification of the abstract Ext basis with the project's fixed E2 product remains the registered E2Computation obligation; this source check does not prove it. The finite archive bound t<=261 is not extrapolated. Primary text: [Adams](https://www.sas.rochester.edu/mth/sites/doug-ravenel/otherpapers/Adams-HI1.pdf), Theorem 2.5.1(2), p.45.
\end{theorem}
